\section{Configuration du coupleur \texttt{NOC\_ETHIP\_CONVOYEUR}}
\label{secConfigurationNocConvoyeur}

Il s'agit maintenant d'appliquer, sous Control Expert, la procédure
générique de configuration d'un coupleur \texttt{BMX~NOC~0401.2}
(document \texttt{02c\_configuration\_noc}) au cas particulier du
coupleur \texttt{NOC\_ETHIP\_CONVOYEUR}, avec les paramètres numériques
propres à cet exercice.

\begin{UPSTIManipulation}{Configuration matérielle et réseau}
    \UPSTIetape{\textbf{Étape~1} (\ref{secEtapeImportEds}) : vérifier que
    le fichier EDS du PFC200 (\emph{WAGO~2759-101~EN/IP}) est bien
    disponible dans le catalogue DTM.}

    \UPSTIetape{\textbf{Étape~2} (\ref{secEtapeConfigMaterielle}) :
    déclarer le coupleur \texttt{BMX~NOC~0401.2} correspondant, en le
    nommant \texttt{NOC\_ETHIP\_CONVOYEUR}.}

    \UPSTIetape{\textbf{Étape~3} (\ref{secEtapeConfigReseau}) : paramétrer
    la voie \texttt{TCP/IP} du coupleur avec les valeurs déjà données pour
    ce réseau (adresse IP, masque, passerelle).}
\end{UPSTIManipulation}

\begin{UPSTIManipulation}{Ajout et paramétrage du PFC200}
    \UPSTIetape{\textbf{Étape~4} (\ref{secEtapeAjoutEquipement}) :
    ajouter l'équipement PFC200 sur le réseau, le nommer \texttt{guichet},
    puis paramétrer le service d'échange \emph{Exclusive Owner} avec une
    taille d'entrée de \textbf{12~octets} et une taille de sortie de
    \textbf{4~octets}, conformément à la volumétrie donnée en
    \ref{secDonneesGuichet}.}

    \UPSTIetape{Renseigner l'adresse IP réelle du PFC200 dans l'onglet
    \emph{Paramétrage de l'adresse}.}

    \UPSTIetape{Choisir le mode d'importation manuel des items, puis nommer
    chacun des bits produits et consommés en réutilisant les références
    données en \ref{secDonneesGuichet} (préfixes \texttt{m\_nix} et
    \texttt{m\_nqx}).}
\end{UPSTIManipulation}

\begin{UPSTIManipulation}{Vérification avant téléchargement}
    \UPSTIetape{\textbf{Étape~5} (\ref{secEtapeVerification}) :
    sauvegarder et compiler l'application.}

    \UPSTIetape{À l'aide de l'éditeur de données, contrôler que les
    variables \texttt{guichet\_IN} et \texttt{guichet\_OUT} sont bien
    positionnées aux adresses \texttt{\%MW} déterminées en
    \cref{secCartographieMemoire}.}

    \UPSTIetape{Faire vérifier la configuration avant tout téléchargement
    dans l'automate.}
\end{UPSTIManipulation}
